\documentclass[10pt,a4paper]{article}
\usepackage[utf8]{inputenc}
\usepackage[T1]{fontenc}
\usepackage{lmodern}
\usepackage[margin=1in]{geometry}
\usepackage{xcolor}
\usepackage{amsmath,amssymb,amsfonts}
\usepackage{graphicx}
\usepackage{booktabs}
\usepackage{adjustbox}
\usepackage{tcolorbox}
\usepackage{enumitem}
\usepackage{microtype}
\usepackage{array}
\usepackage{fancyhdr}

% Custom Colors matching house style
\definecolor{bitsnavy}{RGB}{33,53,94}
\definecolor{bitsblue}{RGB}{44,90,160}
\definecolor{bitsgreen}{RGB}{46,125,82}
\definecolor{bitsamber}{RGB}{200,136,30}
\definecolor{bitsred}{RGB}{178,58,72}

% Callout Boxes
\tcolorboxenvironment{intuition}{colback=bitsblue!8,colframe=bitsblue,title=\textbf{Everyday Picture \& Intuition},fonttitle=\bfseries,boxrule=0.8pt,arc=3pt}
\tcolorboxenvironment{worked}{colback=bitsgreen!8,colframe=bitsgreen,title=\textbf{Fully Worked Example},fonttitle=\bfseries,boxrule=0.8pt,arc=3pt}
\tcolorboxenvironment{everyday}{colback=bitsnavy!5,colframe=bitsnavy,title=\textbf{Real-World Picture},fonttitle=\bfseries,boxrule=0.8pt,arc=3pt}
\tcolorboxenvironment{watchout}{colback=bitsred!8,colframe=bitsred,title=\textbf{Watch Out! Pitfalls},fonttitle=\bfseries,boxrule=0.8pt,arc=3pt}
\tcolorboxenvironment{keytake}{colback=bitsamber!10,colframe=bitsamber,title=\textbf{Key Takeaways},fonttitle=\bfseries,boxrule=0.8pt,arc=3pt}

\newenvironment{intuition}{\begin{tcolorbox}}{\end{tcolorbox}}
\newenvironment{worked}{\begin{tcolorbox}}{\end{tcolorbox}}
\newenvironment{everyday}{\begin{tcolorbox}}{\end{tcolorbox}}
\newenvironment{watchout}{\begin{tcolorbox}}{\end{tcolorbox}}
\newenvironment{keytake}{\begin{tcolorbox}}{\end{tcolorbox}}

\newenvironment{steps}{\begin{enumerate}[label=\textbf{Step \arabic*:},leftmargin=*]}{\end{enumerate}}

\newcommand{\secsub}[1]{\noindent\textit{\large #1}\vspace{0.8em}}
\newcommand{\housefig}[2]{
\begin{figure}[htbp]
\centering
\includegraphics[width=0.85\linewidth]{#1}
\caption{#2}
\end{figure}
}

\pagestyle{fancy}
\fancyhf{}
\rhead{\color{bitsnavy}\bfseries WILP, BITS Pilani}
\lhead{\color{bitsnavy}\bfseries ISM Session 14 Study Companion}
\rfoot{Page \thepage}

\begin{document}

% Banner
\begin{tcolorbox}[colback=bitsnavy,colframe=bitsnavy,arc=0pt,outer arc=0pt,top=12pt,bottom=12pt]
\color{white}
\centering
{\large\bfseries WILP, BITS Pilani}\\[4pt]
{\LARGE\bfseries Session 14: Advanced Time Series (ARIMA, SARIMAX, VAR)}\\[4pt]
{\small Course: Introduction to Statistical Methods (ISM) \quad|\quad Instructor: Dr. YVK Ravi Kumar}
\end{tcolorbox}

\vspace{1em}

\section{White Noise, Dickey-Fuller Test \& ARMA Models}
\secsub{Diagnose white noise residuals, test stationarity, and combine past memory with shocks.}

A time series is white noise if $\epsilon_t \sim \text{i.i.d.}(0, \sigma^2)$ with zero autocovariance. The Dickey-Fuller test checks for a unit root ($\phi = 1$) in $\Delta y_t = (\phi - 1) y_{t-1} + \epsilon_t$. The ARMA($p, q$) model combines autoregressive memory with moving average shock corrections:
\begin{equation}
y_t = c + \phi_1 y_{t-1} + \dots + \phi_p y_{t-p} + \epsilon_t + \theta_1 \epsilon_{t-1} + \dots + \theta_q \epsilon_{t-q}
\end{equation}

\begin{intuition}
White noise is pure static on a radio. The Dickey-Fuller test checks whether your radio station drift is stationary, while ARMA models memory plus lingering surprise shocks.
\end{intuition}

\housefig{figures/acf_pacf_patterns.png}{Autocorrelation Function (ACF) exponential decay vs Partial Autocorrelation (PACF) cutoff.}

\begin{worked}
Given an ARMA(1,1) model: $y_t = 3 + 0.5 y_{t-1} + \epsilon_t + 0.4 \epsilon_{t-1}$.
Last observed value $y_{100} = 10$, and last residual $\epsilon_{100} = 2$.
\begin{steps}
\item Forecast 1 step ahead ($t = 101$):
\begin{equation}
\hat{y}_{101} = 3 + 0.5(10) + 0.4(2) = 3 + 5 + 0.8 = 8.8
\end{equation}
\item Forecast 2 steps ahead ($t = 102$):
Since $\epsilon_{101}$ is unknown, its expected value $E[\epsilon_{101}] = 0$. The MA term drops out:
\begin{equation}
\hat{y}_{102} = 3 + 0.5(8.8) = 3 + 4.4 = 7.4
\end{equation}
Notice how the MA shock term only affects the immediate 1-step forecast!
\end{steps}
\end{worked}

\section{ARIMA \& SARIMAX: Differencing, Seasonality, and External Factors}
\secsub{Stationarize trending data via differencing and layer seasonal and exogenous terms.}

ARIMA($p,d,q$) differences non-stationary series $d$ times ($w_t = \Delta^d y_t$) before fitting ARMA. SARIMAX extends this by adding seasonal periods $m$ and exogenous regressors $X_t$:
\begin{equation}
\Phi_P(B^m) \phi_p(B) \Delta^d \Delta_m^D y_t = \beta X_t + \Theta_Q(B^m) \theta_q(B) \epsilon_t
\end{equation}

\begin{intuition}
ARIMA strips away constant upward trends. SARIMAX adds seasonal calendars (e.g. December holiday surges) and external drivers like promotional marketing discounts.
\end{intuition}

\housefig{figures/arima_sarimax_forecast.png}{Multi-step forecasts comparing ARIMA baseline vs SARIMAX promotion-adjusted forecasts.}

\begin{worked}
Consider an ARIMA(1,1,1) model fitted on differenced series $w_t = y_t - y_{t-1}$:
Fitted equation: $w_t = 0.4 w_{t-1} + \epsilon_t + 0.3 \epsilon_{t-1}$.
Given $y_{99} = 115, y_{100} = 120 \implies w_{100} = 5$. Last error $\epsilon_{100} = 1$.
\begin{steps}
\item Forecast differenced value $\hat{w}_{101}$:
\begin{equation}
\hat{w}_{101} = 0.4(5) + 0.3(1) = 2.0 + 0.3 = 2.3
\end{equation}
\item Un-difference to get level forecast $\hat{y}_{101}$:
\begin{equation}
\hat{y}_{101} = y_{100} + \hat{w}_{101} = 120 + 2.3 = 122.3
\end{equation}
\item Forecast step 2 differenced value $\hat{w}_{102}$ ($E[\epsilon_{101}] = 0$):
\begin{equation}
\hat{w}_{102} = 0.4(2.3) = 0.92 \implies \hat{y}_{102} = 122.3 + 0.92 = 123.22
\end{equation}
\end{steps}
\end{worked}

\section{Vector Autoregression (VAR \& VARMAX)}
\secsub{Model multiple interdependent time series with mutual cross-lag feedback.}

VAR($p$) models a vector of time series $\mathbf{Y}_t$, where each series depends on past values of itself and past values of every other series in the system:
\begin{equation}
\mathbf{Y}_t = \mathbf{c} + \mathbf{A}_1 \mathbf{Y}_{t-1} + \dots + \mathbf{A}_p \mathbf{Y}_{t-p} + \boldsymbol{\epsilon}_t
\end{equation}

\begin{intuition}
GDP growth affects inflation, but inflation also feeds back into GDP growth. VAR treats both as equal partners in a joint dynamic ecosystem.
\end{intuition}

\housefig{figures/var_system_forecast.png}{Multivariate VAR(1) joint forecast tracking interdependencies between GDP growth and inflation.}

\begin{worked}
Given a 2-variable VAR(1) system for GDP growth ($g_t$) and Inflation ($\pi_t$):
\begin{align*}
g_t &= 0.3 + 0.5 g_{t-1} + 0.1 \pi_{t-1} \\
\pi_t &= 0.2 + 0.2 g_{t-1} + 0.6 \pi_{t-1}
\end{align*}
Given last quarter values $g_{99} = 2.0\%$ and $\pi_{99} = 3.0\%$.
\begin{steps}
\item Forecast GDP growth for quarter 100:
\begin{equation}
\hat{g}_{100} = 0.3 + 0.5(2.0) + 0.1(3.0) = 0.3 + 1.0 + 0.3 = 1.6\%
\end{equation}
\item Forecast Inflation for quarter 100:
\begin{equation}
\hat{\pi}_{100} = 0.2 + 0.2(2.0) + 0.6(3.0) = 0.2 + 0.4 + 1.8 = 2.4\%
\end{equation}
\end{steps}
\end{worked}

\begin{watchout}
VAR models require all target series to be stationary. Use the Dickey-Fuller test to verify stationarity before fitting VAR.
\end{watchout}

\begin{keytake}
\begin{itemize}
\item White Noise residuals indicate well-fitted model dynamics.
\item Dickey-Fuller tests verify stationarity.
\item ARMA combines past memory (AR) and shock residuals (MA).
\item ARIMA uses differencing $d$ to stationarize trends; SARIMAX adds seasonal periods $m$ and external regressors $X$.
\item VAR models multivariate time series systems with mutual feedback across variables.
\end{itemize}
\end{keytake}

\end{document}
